% =========================================================================
% Figura: Soluciones de Pareto y soluciones dominadas en el espacio de objetivos
% =========================================================================

\begin{figure}[H]
\centering
\resizebox{0.60\linewidth}{!}{%
\begin{tikzpicture}[
    font=\sffamily,
    >=Stealth,
    x=1cm, y=1cm
]

% -------------------------------------------------------------------------
% 1. EJES COORDENADOS (Con indicación explícita de minimización)
% -------------------------------------------------------------------------
% Eje horizontal f1
\draw[line width=1.1pt, draw=darkText, ->, -{Stealth[length=6pt, width=4pt]}]
    (0, 0) -- (8.8, 0);
\node[below=8pt, font=\fontsize{10}{12}\selectfont, text=darkText]
    at (4.4, 0) {\textbf{Objetivo $f_1$} (minimizar)};

% Eje vertical f2
\draw[line width=1.1pt, draw=darkText, ->, -{Stealth[length=6pt, width=4pt]}]
    (0, 0) -- (0, 6.6);
\node[font=\fontsize{10}{12}\selectfont, text=darkText, rotate=90]
    at (-0.65, 3.3) {\textbf{Objetivo $f_2$} (minimizar)};

% -------------------------------------------------------------------------
% 2. LEYENDA (Frente de Pareto y Soluciones dominadas)
% -------------------------------------------------------------------------
\begin{scope}[shift={(3.8, 4.6)}]
    \draw[draw=grayText!70, fill=white, rounded corners=5pt, line width=1.0pt]
        (0, 0) rectangle (4.7, 1.45);
    
    % Frente de Pareto (Línea + punto azul)
    \draw[line width=1.1pt, draw=figBlue] (0.20, 1.0) -- (0.60, 1.0);
    \filldraw[draw=figBlue!90!black, fill=figBlue, line width=0.8pt]
        (0.40, 1.0) circle (0.12cm);
    \node[anchor=west, text=darkText, font=\fontsize{9}{11}\selectfont]
        at (0.72, 1.0) {\textbf{Frente de Pareto}};
    
    % Soluciones dominadas (Punto naranja)
    \filldraw[draw=figOrange!90!black, fill=figOrange, line width=0.8pt]
        (0.40, 0.45) circle (0.12cm);
    \node[anchor=west, text=darkText, font=\fontsize{9}{11}\selectfont]
        at (0.72, 0.45) {\textbf{Soluciones dominadas}};
\end{scope}

% -------------------------------------------------------------------------
% 3. SOLUCIONES DOMINADAS (Bolitas color naranja)
% -------------------------------------------------------------------------
\newcommand{\drawDominated}[2]{%
    \filldraw[draw=figOrange!90!black, fill=figOrange, line width=0.8pt]
        (#1, #2) circle (0.12cm);%
}

\drawDominated{0.65}{5.60}
\drawDominated{1.10}{4.95}
\drawDominated{1.30}{4.30}
\drawDominated{1.75}{4.60}
\drawDominated{2.10}{4.10}
\drawDominated{1.95}{2.25}
\drawDominated{3.00}{1.95}
\drawDominated{3.25}{1.65}
\drawDominated{3.80}{1.40}
\drawDominated{4.55}{1.00}
\drawDominated{5.50}{0.75}
\drawDominated{6.70}{0.50}

% -------------------------------------------------------------------------
% 4. FRENTE DE PARETO (Trazo continuo + Bolitas color azul)
% -------------------------------------------------------------------------
% Línea que une y delimita la frontera no dominada
\draw[draw=figBlue, line width=1.1pt]
    (0.25, 6.00) -- (0.50, 5.30) -- (0.70, 4.80) -- (0.90, 4.50) --
    (1.05, 3.80) -- (1.15, 3.00) -- (1.30, 2.40) -- (1.90, 1.95) --
    (2.75, 1.85) -- (3.05, 1.45) -- (3.75, 1.05) -- (4.35, 0.85) --
    (5.25, 0.60) -- (6.25, 0.40) -- (6.65, 0.32) -- (7.35, 0.20);

\newcommand{\drawPareto}[2]{%
    \filldraw[draw=figBlue!90!black, fill=figBlue, line width=0.8pt]
        (#1, #2) circle (0.12cm);%
}

\drawPareto{0.25}{6.00}
\drawPareto{0.50}{5.30}
\drawPareto{0.70}{4.80}
\drawPareto{0.90}{4.50}
\drawPareto{1.05}{3.80}
\drawPareto{1.15}{3.00}
\drawPareto{1.30}{2.40}
\drawPareto{1.90}{1.95}
\drawPareto{2.75}{1.85}
\drawPareto{3.05}{1.45}
\drawPareto{3.75}{1.05}
\drawPareto{4.35}{0.85}
\drawPareto{5.25}{0.60}
\drawPareto{6.25}{0.40}
\drawPareto{6.65}{0.32}
\drawPareto{7.35}{0.20}

\end{tikzpicture}%
}
\caption{Frente de Pareto (soluciones no dominadas) y soluciones dominadas en el espacio de objetivos ($f_1$ y $f_2$ a minimizar). El frente delimita el límite óptimo de compromiso entre ambos criterios, adaptado de~\cite{talbi2009metaheuristics}.}
\label{fig:frontera_pareto}
\end{figure}
